% Fragmento público generado desde propietarios íntegros.
% Residencia: 52.5.1.
% Fuente: ../incoming/radion_monografia_20260827/manuscrito/sections/06_interior_helice.tex:1-108
\noindent La géométrie projective distingue l'action finie de la profondeur hyperbolique et de la mémoire du retour.\par\medskip

\subsubsection{Intérieur hyperbolique de Hilbert--Beltrami--Klein}

\paragraph{Normalisation projective dans le disque de Klein}

L'application affine
\[
u=\frac{Q}{a_S},\qquad v=\frac{P}{b_S}
\]
envoie l'ellipse d'action sur le disque unité
\[
D=\{(u,v):u^2+v^2<1\}.
\]
Sur chaque corde, la distance de Hilbert est définie par le birapport
de ses extrémités de frontière. La métrique infinitésimale résultante est la forme
de Beltrami--Klein
\begin{equation}
\diff s_K^2=
\frac{(1-r^2)(\diff u^2+\diff v^2)+(u\diff u+v\diff v)^2}
{(1-r^2)^2},
\qquad r^2=u^2+v^2.
\label{eq:klein-metric}
\end{equation}
Avec cette normalisation, la courbure gaussienne est constante :
\begin{equation}
K=-1.
\label{eq:klein-curvature}
\end{equation}
Les géodésiques sont des cordes euclidiennes, mais leur longueur hyperbolique diverge
à l'approche de la frontière.

\paragraph{Finitude de la frontière et infinitude de la profondeur hyperbolique}

Le déterminant métrique de \eqref{eq:klein-metric} produit
\begin{equation}
\diff A_K=\frac{\diff u\,\diff v}{(1-r^2)^{3/2}}.
\label{eq:klein-area-density}
\end{equation}
L'aire du sous-disque \(r<R<1\) est
\begin{align}
A_K(R)
&=\int_0^{2\pi}\int_0^R
\frac{r\,\diff r\,\diff\theta}{(1-r^2)^{3/2}}\\
&=2\pi\left(\frac1{\sqrt{1-R^2}}-1\right).
\label{eq:klein-area-R}
\end{align}
\Needspace{4\baselineskip}
Par conséquent,
\[
\lim_{R\to1^-}A_K(R)=+\infty.
\]
Simultanément, la même frontière enferme l'aire symplectique
\[
\pi a_Sb_S=2\pi\hret=\hfull<\infty.
\]

\begin{theorem}[Finitude de l'action et infinitude de la profondeur]
L'ellipse du radion possède une action de tour finie \(\hfull\), tandis que
son intérieur, muni de la métrique de Hilbert--Beltrami--Klein, a une aire
hyperbolique infinie :
\begin{equation}
\boxed{
\operatorname{Area}_{\mathrm{simp}}(\partial\mathbb E_S)=\hfull<\infty,
\qquad
\operatorname{Area}_{K}(\mathbb E_S)=+\infty.}
\label{eq:finite-infinite}
\end{equation}
\end{theorem}

Il n'y a pas de contradiction : \(\operatorname{Area}_{\mathrm{simp}}\) et
\(\operatorname{Area}_K\) sont des mesures différentes. La première compte l'action
orientée en phase ; la seconde mesure la profondeur projective intérieure. La thèse
\enquote{bord fini, intérieur infini} est désormais une égalité et une limite,
non une métaphore.

\paragraph{Coordonnées focales dans l'intérieur de Klein}

Dans le disque normalisé, les foyers se situent en \((\pm e_f,0)\). La distance
du centre à un foyer est
\begin{equation}
u_f=\operatorname{artanh}e_f
=\operatorname{arcosh}(e^\eta)
=0.813382331175342333760889731088228\ldots.
\label{eq:klein-focal}
\end{equation}
La distance foyer--foyer est
\[
d_K(F_-,F_+)=2u_f
=1.62676466235068466752177946217646\ldots.
\]
On a
\begin{equation}
e_f=\tanh u_f,\qquad
e^{-\eta}=\operatorname{sech}u_f,\qquad
\eta=\log\cosh u_f.
\label{eq:focal-hyperbolic-chain}
\end{equation}
La région jusqu'au rayon focal a pour aire
\[
A_K(e_f)=2\pi(e^\eta-1)
=2.19559620884061869541206115980360\ldots.
\]

Toutes les ellipses ouvertes sont projectivement isométriques en tant que domaines de
Hilbert. C'est pourquoi la courbure \(-1\) seule ne mémorise pas l'excentricité. La
forme HMT réside dans la structure conjointe : carte affine distinguée, forme
symplectique, étiquetage des feuillets et paramétrisation de frontière.
